\documentclass[10pt,a4paper]{amsart}
\usepackage[margin=19mm]{geometry}
\usepackage{amsmath,amssymb,mathtools}
\usepackage[scr=rsfs,cal=euler]{mathalfa}
\usepackage{xfrac}
\usepackage[unicode,hidelinks,bookmarks=false]{hyperref}
\newcommand{\ie}{i.e.\ }
\newcommand{\cf}{cf.\ }
\newcommand{\resp}{respectively\ }
\newcommand{\Subsection}{\subsection}
\providecommand{\nobreakdash}{\nobreak\hbox{-}}
\setlength{\emergencystretch}{2em}
\multlinegap0pt
\usepackage{bm}
\usepackage{bbm}
\usepackage{dsfont}
\usepackage{xspace}
\usepackage{cancel}
\usepackage{mathtools}
\usepackage{amsthm}
\makeatletter
\@ifundefined{theorem}{\newtheorem{theorem}{Theorem}}{}
\@ifundefined{lemma}{\newtheorem{lemma}[theorem]{Lemma}}{}
\makeatother
\title[Thurston spine in a Teichm\"uller curve]{Thurston spine in a Teichm\"uller curve}
\author{Yue Gao, Zhongzi Wang}
\date{Source: arXiv:2509.24141v1 (CC BY 4.0)}
\begin{document}
\maketitle

\noindent\emph{Source and licence.} Thurston spine in a Teichm\"uller curve. Version arXiv:2509.24141v1 (CC BY 4.0), distributed under the Creative Commons Attribution 4.0 International licence, as recorded in the arXiv licence metadata for this exact version and shown on the abstract page https://arxiv.org/abs/2509.24141v1. Full rights notice: licenses/arxiv-2509.24141-CC-BY-4.0.txt.\par\smallskip

\noindent\textit{Editorial scope.} Counted unit: the Lemma whose source label is \texttt{lem:limit} in the article (source0.tex lines 507--510, source label \texttt{lem:limit}), delivered together with the article's own complete original proof of it (source0.tex lines 511--582, one \texttt{proof} environment of 72 lines). No mathematics is written, repaired or re-derived: statement and proof are the article's own text line for line. Required context retained uncounted: thm:minsky (lines 422--433). Every definition, display and label of those lines is retained unchanged. Omitted from this package and not redistributed: 1--421, 434--506, 583--1589. Nothing inside a retained excerpt is omitted or altered: every retained body is delivered exactly as the article's lines are written, so no omission receipt of any kind accompanies this package. Citations: the retained excerpts carry no citation command at all, so no bibliography entry of the article is transcribed and no reference is redistributed. Reference adaptation: none was needed and none was made; every reference inside the delivered unit resolves inside it, so no unresolved reference or citation is produced. Printed slips kept literal: no printed word, symbol, hypothesis or number of the counted statement or of its proof was altered. Layout and class: the article's own class is \texttt{amsart} with packages including amsmath, amssymb, graphicx, bbm, amsthm, caption, subcaption, babel, color, hyperref, xy, tikz-cd, xcolor, verbatim; the wrapper is the lane's pinned \texttt{amsart} editorial wrapper at 10pt on A4, which restates the theorem-like environments the retained text needs and adds the article's own macro definitions that the retained text needs (none), with extra wrapper packages (bm, bbm, dsfont, xspace, cancel, mathtools). The delivered numbering is the unit's own. Assets: no retained span contains a figure environment, a graphics-inclusion command, a PDF-page inclusion, a table or a TikZ picture; no image file is copied into this package, the package declares no asset dependency and no image file is redistributed with it. Licence: version arXiv:2509.24141v1 of this article is distributed under the Creative Commons Attribution 4.0 International licence, as recorded in the arXiv licence metadata for this exact version and shown on the abstract page https://arxiv.org/abs/2509.24141v1. A full rights notice is delivered at licenses/arxiv-2509.24141-CC-BY-4.0.txt. Characters: every retained excerpt is pure ASCII, so the delivered engine, XeLaTeX with its default Latin Modern text font, reports no missing character.
\begin{theorem}
 For the Fenchel-Nielsen coordinate $(c,t)$ of $\mathcal{T}_{0,4}$, identify $\mathcal{T}_{0,4}$ and $\mathbb{H}^2$ by 
 \begin{eqnarray*}
  \phi:\mathcal{T}_{0,4} & \to & \mathbb{H}^2 \\
  (c,t) &\mapsto&   \frac{t}{c} + \frac{i}{c}. 
 \end{eqnarray*}
 Let $d_{\mathcal{T}}$ be the Teichm\"uller distance on $\mathcal{T}_{0,4}$ and $d_{\mathbb{H}}$ be the hyperbolic distance on $\mathbb{H}^2$. For a sufficiently small $\epsilon >0$, for any $X, Y \in \mathcal{T}_{0,4}$ satisfying $\ell_{\gamma}(X), \ell_{\gamma}(Y) < \epsilon$, there is a constant $C >0$ such that
 \[
 |d_{\mathcal{T}}(X,Y) - d_{\mathbb{H}}(\phi(X), \phi(Y))| < C. 
 \]
 \label{thm:minsky}
\end{theorem}

\begin{lemma}
 The geodesic rays $L_{-1}|_{c \le 1}$, $L_0|_{c \le 1}$, $L_1|_{c \le 1}$ are the same point of the Gromov boundary of $\mathcal{T}_{0,4}$. 
 \label{lem:limit}
\end{lemma}

\begin{proof}
 The main ingredient of the proof is Theorem \ref{thm:minsky}. Some symbols for example $\phi$, $C$, $d_{\mathcal{T}}$, $d_{\mathbb{H}}$ are from the statement of Theorem \ref{thm:minsky}. 
Take the Fenchel-Nielsen coordinate $(c,t)$ of $\mathcal{T}_{0,4}$.
 For a sufficiently small $c_0 \in (0,1)$,  let
 \begin{eqnarray*}
  l_j: (0,c_0] &\to& \mathcal{T}_{0,4} \\
  c &\mapsto& (c,j\cdot c)
 \end{eqnarray*}
 be the map corresponding to the concerned geodesic ray $L_j|_{c \le c_0}$, where $j= -1, 0, 1$. 
 Let 
 \begin{eqnarray*}
  f_j: (0,c_0] & \to & [0,+\infty) \\
  c& \mapsto & x = f_j(c) 
 \end{eqnarray*}
 be the arc length reparametrization of $l_j$. Namely
 \begin{eqnarray*}
  l_j\circ f_j^{-1}: [0,+\infty) & \to & \mathcal{T}_{0,4} \\
         x & \mapsto &  l_j\circ f_j^{-1}(x)= l_j(c) = (c, j\cdot c)
 \end{eqnarray*}
 is the geodesic ray parametrized by arc length with respect to the Teichm\"uller metric $d_{\mathcal{T}}$. 

 WLOG, one may consider the rays $L_0|_{c \le c_0}$ and $L_1|_{c \le c_0}$. 
 The endpoints of these rays in $\mathcal{T}_{0,4}$ are $$l_0\circ f^{-1}_0(0) = (c_0, 0)\text{ and }l_1\circ f^{-1}_1(0) = (c_0, 1\cdot c_0)$$ respectively. For these points, one may have $$\phi\circ l_0\circ f^{-1}_0(0) = \frac{i}{c_0}\text{ and }\phi\circ l_1\circ f^{-1}_1(0) = 1+\frac{i}{c_0}.$$ 
 For any $x>0$, consider $$l_0\circ f^{-1}_0(x) = (f_0^{-1}(x),0) \text{ and }l_1\circ f^{-1}_1(x) = (f_1^{-1}(x),1\cdot f_1^{-1}(x)) .$$
 For these points, one may have that
\begin{equation}
\phi\circ l_0\circ f^{-1}_0(x) =  \frac{i}{f_0^{-1}(x)} \text{ and }\phi\circ l_1\circ f^{-1}_1(x) = 1+ \frac{i}{f_1^{-1}(x)}. 
\label{for:hcoord}
\end{equation}
 Since $c_0$ is sufficiently small, by Theorem \ref{thm:minsky}, there is a $C>0$ such that, on the geodesic ray $L_0|_{c \le c_0}$, one may have that
 \[
  |d_{\mathcal{T}}(l_0\circ f^{-1}_0(x), l_0\circ f^{-1}_0(0)) - d_{\mathbb{H}} (\phi\circ(l_0\circ f^{-1}_0)(x), \phi\circ(l_0\circ f^{-1}_0)(0)) |  < C 
 \]
namely,
 \[
 \left|x - \log \frac{c_0}{f_0^{-1}(x)} \right| <C. 
 \]
On the geodesic ray $L_1|_{c \le c_0}$, one may have that
 \[
 |d_{\mathcal{T}}(l_1\circ f^{-1}_1(x), l_1\circ f^{-1}_1(0)) - d_{\mathbb{H}} (\phi\circ(l_1\circ f^{-1}_1)(x), \phi\circ(l_1\circ f^{-1}_1)(0)) |  < C 
 \]
 namely,
 \[
 \left|x - \log \frac{c_0}{f_1^{-1}(x)} \right| <C.   
 \]
 Eliminating $c_0$ and $x$, one may get that
 \begin{equation}
  \left| \log \frac{f^{-1}_1(x)}{f^{-1}_0(x)} \right| < 2C . 
  \label{for:lipshitz}
 \end{equation}
 Again by Theorem \ref{thm:minsky}, one may have that
 \begin{equation}
  |d_{\mathcal{T}}(l_0\circ f_0^{-1}(x), l_1\circ f_1^{-1}(x)) - d_{\mathbb{H}}(\phi\circ l_0\circ f_0^{-1}(x), \phi\circ l_1\circ f_1^{-1}(x))| < C. 
  \label{for:dtdh}
 \end{equation}
 Then, one may get that
 \begin{eqnarray}
  & & d_{\mathbb{H}}(\phi\circ l_0 \circ f_0^{-1}(x), \phi\circ l_1 \circ f_1^{-1}(x))\nonumber \\
  &=& d_{\mathbb{H}}\left( \frac{i}{f_0^{-1}(x)}, 1+\frac{i}{f_1^{-1}(x)} \right) \nonumber\,\,\,\,\text{(by (\ref{for:hcoord}))}\\
  & \le & d_{\mathbb{H}}\left( \frac{i}{f_0^{-1}(x)}, 1+\frac{i}{f_0^{-1}(x)} \right) + d_{\mathbb{H}}\left( 1+\frac{i}{f_0^{-1}(x)}, 1+\frac{i}{f_1^{-1}(x)} \right) \nonumber\\
  & \le & 1+2C. \,\,\,\,\text{(by }f_0^{-1}(x) < c_0 <1 \text{ and (\ref{for:lipshitz}) )} \label{for:dh}
 \end{eqnarray}
 Combining (\ref{for:dtdh}) and (\ref{for:dh}), one may get that, for any $x>0$, 
 \[
 d_{\mathcal{T}}(l_0\circ f_0^{-1}(x), l_1\circ f_1^{-1}(x)) \le 1+3C.  
 \]

 Therefore the geodesic rays $L_{0}|_{c \le 1}$ and $L_{1}|_{c \le 1}$ represent the same point in the Gromov boundary. 
The same proof holds for $L_{-1}|_{c \le 1}$ and  $L_0|_{c \le 1}$. The proof is complete. 


\end{proof}

\end{document}
